\documentclass[english]{article}
\usepackage[T1]{fontenc}
\usepackage[utf8]{inputenc}
\setcounter{secnumdepth}{-2}
\setlength{\parindent}{0bp}
\usepackage{amsmath}
\usepackage{amsthm}
\usepackage{amssymb}

\makeatletter
\theoremstyle{plain}
\newtheorem*{thm*}{\protect\theoremname}

\@ifundefined{date}{}{\date{}}
\usepackage{graphicx}
\usepackage{geometry}
\newcommand{\limi}{\underset{n\rightarrow\infty}{\lim}}
\newcommand{\limxz}{\underset{x\rightarrow x_{0}}{\lim}}
\newcommand{\limfxz}{\underset{x\rightarrow x_{0}}{\lim}f(x)}
\newcommand{\limfxm}{\underset{x\rightarrow x_{0}^{-}}{\lim}f(x)}
\newcommand{\limfxp}{\underset{x\rightarrow x_{0}^{+}}{\lim}f(x)}
\newcommand{\limxm}{\underset{x\rightarrow x_{0}^{-}}{\lim}}
\newcommand{\limxp}{\underset{x\rightarrow x_{0}^{+}}{\lim}}
\newcommand{\sqrm}{M_{m\times m}(\mathbb{F})}
\newcommand{\seqa}{(a_{n})_{n=1}^{\infty}}
\newcommand{\seqx}{(x_{n})_{n=1}^{\infty}}
\newcommand{\funcdr}{f:D\rightarrow\mathbb{R}}
\newcommand{\der}{\limxz\frac{f(x)-f(x_{0})}{x-x_{0}}}
\usepackage[all]{pdfprivacy}

\makeatother

\usepackage{babel}
\providecommand{\theoremname}{Theorem}

\begin{document}
\newgeometry{left=2cm, right=2cm, top=2cm, bottom=2cm}
\title{{\Large\vspace{-4em}The intermediate value theorem\vspace{-2em}}}
\maketitle
\begin{thm*}
Let $a,b,\gamma\in\mathbb{R}$ s.t. $a<b$ and let $f:D\rightarrow\mathbb{R}$
be a continuous function s.t. $\left[a,b\right]\subseteq D$ and $f(a)\ne f(b)$. 

If $f(a)<\gamma<f(b)$ or $f(b)<\gamma<f(a)$ then there exists $c\in(a,b)\text{ s.t. }f(c)=\gamma$.
\end{thm*}
\begin{proof}
Let us define $g:D\rightarrow\mathbb{R}$ by $g(x)=f(x)-\gamma$.

We will first prove that $\exists x_{0}\in\left(a,b\right)\text{ s.t. }g\left(x_{0}\right)=0$.

$g$ is continuous as a shift of a continuous function. 

Let us assume that $f(a)<\gamma<f(b)$. Then
\[
g(a)=f(a)-\gamma<0<f(b)-\gamma=g(b)
\]
Let us define the two sequences $\left(a_{n}\right)_{n=1}^{\infty}$
and $\left(b_{n}\right)_{n=1}^{\infty}$ inductively. 

$a_{1}=a,b_{1}=b$.

For $n\in\mathbb{N}$,

We will assume that $g\left(\frac{a_{n}+b_{n}}{2}\right)\ne0$, otherwise
$x_{0}=\frac{a_{n}+b_{n}}{2}$ and the theorem holds.

If $g\left(\frac{a_{n}+b_{n}}{2}\right)<0$ we will define $a_{n+1}=\frac{a_{n}+b_{n}}{2}$
and $b_{n+1}=b_{n}$, otherwise $g\left(\frac{a_{n}+b_{n}}{2}\right)>0$
and we will define $a_{n+1}=a_{n}$ and $b_{n+1}=\frac{a_{n}+b_{n}}{2}$.

By induction, the following properties are true for each $n\in\mathbb{N}$:

1. $a_{n}\le a_{n+1}$

2. $b_{n}\ge b_{n+1}$

3. $g\left(a_{n}\right)<0<g\left(b_{n}\right)$

4. $a\le a_{n}<b_{n}\le b$ hence $b_{n}-a_{n}>0$.

We will now look at the sequence $\left(d_{n}\right)_{n=1}^{\infty}$
where $d_{n}=b_{n}-a_{n}$.

Let us notice that in the first case
\[
b_{n+1}-a_{n+1}=b_{n}-\frac{a_{n}+b_{n}}{2}=\frac{2b_{n}-a_{n}-b_{n}}{2}=\frac{b_{n}-a_{n}}{2}
\]

and in the second case
\[
b_{n+1}-a_{n+1}=\frac{a_{n}+b_{n}}{2}-a_{n}=\frac{a_{n}+b_{n}-2a_{n}}{2}=\frac{b_{n}-a_{n}}{2}
\]

So 
\[
\forall n\in\mathbb{N},\ d_{n+1}=\frac{d_{n}}{2}
\]

and $\left(d_{n}\right)_{n=1}^{\infty}$ is a geometric sequence with
$p=\frac{1}{2}$ and $d_{1}=b-a$, hence
\[
\limi d_{n}=\limi\frac{b-a}{2^{n-1}}=0
\]

Points 1 and 2 and 4 give
\[
a_{n}\le a_{n+1}\le b_{n+1}\le b_{n}
\]
and along with what we just proved, Cantor's intersection theorem
provides a $x_{0}\in\left[a,b\right]$ such that
\[
x_{0}=\limi a_{n}=\limi b_{n}
\]
Therefore, according to Heine's sequence characterisation of continuity,
\[
g(x_{0})=\limi g(a_{n})\le0\ \wedge\ g(x_{0})=\limi g(b_{n})\ge0
\]
therefore 
\[
g(x_{0})=0
\]

and because $g(a)<0<g(b)$, $x_{0}\in\left(a,b\right)$.

According to the definition of $g$,
\[
f(x_{0})-\gamma=0\Rightarrow f(x_{0})=\gamma
\]

as required.

If $f(b)<\gamma<f(a)$, we can apply the same proof for $-f$ and
$-\gamma$ because $-f$ is also continuous and 
\[
-f(a)<-\gamma<-f(b)
\]
\end{proof}

\end{document}
